\chapter{The memory cut-off time}
\label{ch:cutoff}

A \emph{memory cut-off time} is a threshold after which a stored EPR pair is discarded and a new pair generated. Nothing new needs to be built to have one: ACE already has it, as the SeQUeNCe memory's
\code{cutoff_ratio}, and two other timers also end a stored pair's life.

\section{The timers}
\begin{description}
  \item[Memory cut-off.] When a pair is stored, an expiry is scheduled \code{cutoff_ratio} $\times$ coherence time later. At expiry the memory is reset (unless it is in an application) and returned to the free state; ACE then drops the
        pair from its cache and tells the compiler controller it expired. ACE's configuration never sets the ratio, so it is $1$ and the cut-off equals the coherence time, $1$\,s. A ratio written in the topology template was verified to reach every memory.
  \item[Reservation window.] A pair's booking ends with its window; an unused pair is released then. The window is $1000$\,ms for compiler pairs and, for the autonomous caches, runs to the next $0.2$\,s period boundary.
  \item[Planner lifetime.] The planner's \code{--coherence-time-layers} (10) says how far ahead a pair may be prepared; it is a planning assumption, tied to nothing physical.
\end{description}
Separately the stored state decoheres continuously at the rate $1/$coherence time, lowering fidelity but discarding nothing.

\section{What a discard does to compiler pairs}
\begin{table}[H]
\centering\small
\resizebox{\ifdim\width>\linewidth\linewidth\else\width\fi}{!}{%
\begin{tabular}{@{}llrrrrrrr@{}}
\toprule
Workload & Case & Done & Viol. & Consumed & Window & Cut-off & Bookings & Makespan ms \\
\midrule
ghz & ACE default: cut-off = coherence time (1 s), window 1000 ms & 24/24 & 0 & 48 & 0 & 0 & 48 & 303.5 \\
ghz & short reservation window only (20 ms) & 24/24 & 0 & 36 & 12 & 0 & 48 & 311.2 \\
ghz & cut-off ratio 0.02 (20 ms), window 1000 ms & 24/24 & 0 & 34 & 12 & 12 & 55 & 311.4 \\
ghz & cut-off ratio 0.005 (5 ms), window 1000 ms & 24/24 & 0 & 0 & 34 & 34 & 49 & 322.0 \\
ghz & cut-off ratio 0.02 and window 20 ms & 24/24 & 0 & 36 & 12 & 0 & 48 & 311.2 \\
graphstate & ACE default: cut-off = coherence time (1 s), window 1000 ms & 76/76 & 0 & 104 & 0 & 0 & 134 & 274.5 \\
graphstate & short reservation window only (20 ms) & 76/76 & 0 & 48 & 56 & 0 & 134 & 241.3 \\
graphstate & cut-off ratio 0.02 (20 ms), window 1000 ms & 76/76 & 0 & 24 & 32 & 32 & 179 & 351.8 \\
graphstate & cut-off ratio 0.005 (5 ms), window 1000 ms & 76/76 & 0 & 0 & 34 & 34 & 339 & 347.4 \\
graphstate & cut-off ratio 0.02 and window 20 ms & 76/76 & 0 & 48 & 56 & 0 & 134 & 242.5 \\
\bottomrule
\end{tabular}}
\caption{Which timer discards a cached pair (compiler pre-generation, 4 memories). Consumed: pairs used; Window: releases at the end of the reservation window; Cut-off: memory resets by the cut-off.}
\label{tab:cutoff-mechanism}
\end{table}

The earlier timer wins: with both at $20$\,ms the window fires first and no cut-off reset occurs. For compiler pairs a cut-off discards the pair and nothing replaces it: the controller only marks it expired, the request that wanted it
generates its own pair on demand, and all requests still complete. The run gets longer (graph state: $274.5$\,ms makespan by default, $351.8$\,ms with a $20$\,ms cut-off, $347.4$\,ms with $5$\,ms). In every case the number of window releases equals the number of cut-off resets,
consistent with a discarded pair's booking staying until its window ends (inferred from the counts).

\section{What a discard does to the autonomous caches}
\begin{table}[H]
\centering\small
\resizebox{\ifdim\width>\linewidth\linewidth\else\width\fi}{!}{%
\begin{tabular}{@{}llrrrrrr@{}}
\toprule
Strategy & Cut-off ratio & Generated & Used & Reported waste & Real unused & Latency ms & Fidelity \\
\midrule
continuous & 1 (default) & 44936 & 90 & 99.7\% & 99.80\% & 5.147 & 0.9297 \\
adaptive-continuous & 1 (default) & 5184 & 120 & 97.8\% & 97.69\% & 4.882 & 0.9258 \\
continuous & 0.2 & 44936 & 90 & 99.7\% & 99.80\% & 5.147 & 0.9297 \\
adaptive-continuous & 0.2 & 5184 & 120 & 97.8\% & 97.69\% & 4.882 & 0.9258 \\
continuous & 0.02 & 447388 & 122 & 10.0\% & 99.97\% & 4.865 & 0.9354 \\
adaptive-continuous & 0.02 & 41472 & 122 & 12.0\% & 99.71\% & 4.865 & 0.9354 \\
\bottomrule
\end{tabular}}
\caption{Autonomous caches on Random-16 (4 memories, coherence 1 s): effect of the memory cut-off. Reported waste omits cut-off discards.}
\label{tab:cutoff-cache}
\end{table}

A cut-off of $200$\,ms changes nothing, because the cache's window already ends at the next $0.2$\,s boundary. At $20$\,ms the cache discards and \emph{regenerates}: ten times more pairs are generated, the used pairs are fresher (fidelity $0.9297$ to $0.9354$) and latency
falls slightly ($5.15$ to $4.87$\,ms). ACE's waste counter does not see cut-off discards (it is incremented only when a generation rule expires at the window end), so the reported waste drops from $99.7\%$ to $10.0\%$ although $99.97\%$ of the generated pairs went unused;
waste measured with a cut-off must be computed as $1-\text{used}/\text{generated}$.

\section{Decisions}
Whether ``cut-off'' should mean discard-and-regenerate (the caches already do; compiler pre-generation falls back on demand and re-preparing would be a new feature); exposing the cut-off in milliseconds as a recorded study parameter; tying the planner's lifetime to it
and fixing the waste counter first; and the coherence time, which comes before any cut-off: with $1$\,s coherence a cut-off never acts, and with the group's $1$\,ms a cut-off equal to the coherence time would discard every cached pair before use.
